% Estilos compartidos por los reportes RE_x-y y las presentaciones: paleta, estilos TikZ,
% macros tipográficas y marcas de trabajo pendiente. Requiere xcolor y tikz ya cargados.

\newcommand{\code}[1]{\texttt{#1}}
\newcommand{\file}[1]{\texttt{#1}}

% === Paleta de los esquemas ========================================
\definecolor{azul}{HTML}{1E40AF}
\definecolor{celeste}{HTML}{00B4D8}
\definecolor{verde}{HTML}{22C55E}
\definecolor{rosa}{HTML}{FF5FA2}
\definecolor{ambar}{HTML}{F59E0B}
\definecolor{gris}{HTML}{6B7280}
\tikzset{
    caja/.style={draw=black!60, rounded corners=2pt, fill=white, align=center,
                 minimum height=7mm, inner sep=4pt, font=\small},
    etapa/.style={caja, fill=azul!8, draw=azul!70},
    dato/.style={caja, fill=black!4, draw=black!50, font=\small\ttfamily},
    decision/.style={caja, fill=ambar!15, draw=ambar!80!black},
    descarte/.style={caja, fill=rosa!12, draw=rosa!70!black, font=\footnotesize},
    flecha/.style={-{Stealth[length=2.2mm]}, thick, draw=black!65},
    nota/.style={font=\footnotesize, text=black!70, align=left},
}

% Separador tipográfico (el «·» directo no existe en T1).
\newcommand{\sep}{\,\textperiodcentered\,}

% === Marcas de trabajo pendiente ===================================
% \todocita{...}: dónde hace falta una referencia y para qué. Visibles en el PDF.
\newcommand{\todocita}[1]{\textcolor{red!75!black}{\small[\textsc{cita}: #1]}}
\newcommand{\todo}[1]{\textcolor{orange!80!black}{\small[\textsc{todo}: #1]}}
